%------------------------------------------------------------------------------%
\section{Algumas Funções Usuais}

A seguir serão apresentados algumas funções usuais, onde algumas delas serão
estudadas com mais detalhes em aulas posteriores.

%------------------------------------------------------------------------------%
\subsection{Função Potência}

A função
\[
  f(x)=x^n,
\]
$n$ é um número natural, é chamada de
\emph{função potência}\index{Função!potência}.
A Figura~\ref{pot} mostra o gráfico de algumas funções potências.

%------------------------------------------------------------------------------%
\begin{figure}
  \centering
  \includegraphics[scale=1.6]{pot1}
  \includegraphics[scale=1.6]{pot2}
  \caption{???}
  \label{pot}
\end{figure}
%------------------------------------------------------------------------------%

O domínio de uma função potência é $\R$.

O conjunto imagem é $\R$ se o expoente for ímpar e é
$[0, \infty]$ se o expoente é par.

Há um único zero em $x=0$.

Quando o expoente é par, $f$ é par, quando o expoente é ímpar, $f$ é
ímpar.

Quando o expoente é par, $f$ é decrescente para $x<0$, e crescente para
$x>0$, quando o expoente é ímpar, $f$ é crescente em $(- \infty, \infty)$.

%------------------------------------------------------------------------------%
\subsection{Função Polinomial}

A função
\[
  f(x) = a_nx^n + a_{n-1}x^{n-1} + \cdots+ a_2x^2 + a_1x+a_0,
\]
onde $n$ é um número inteiro não-negativo e
$a_1$, $a_2$, \dots, $a_n$ são constantes reais.

O domínio de $f$ é $\R$.

Pode ter até $n$ zeros.

As funções do primeiro grau e quadráticas são
exemplos de funções polinomiais.

%------------------------------------------------------------------------------%
\subsection{Função do 1\textsubscript{o} Grau}

É uma função do tipo
\[
  f(x) = ax + b,
\]
com $a, b \in \R$.
O gráfico de uma função do 1\textsuperscript{o} grau é uma reta. $a$ é o
coeficiente angular da reta e $b$ é a intersecção com o eixo $y$. Para
encontrar a equação de uma reta precisamos de dois pontos, em particular se
encontrarmos os interceptos com os eixos conseguimos obter a equação da reta.

%------------------------------------------------------------------------------%
\begin{figure}
  \centering
  \includegraphics[scale=0.35]{linear1} \quad
  \includegraphics[scale=0.35]{linear2}
  \caption{Gráfico de uma função do primeiro grau.}
\end{figure}
%------------------------------------------------------------------------------%

%------------------------------------------------------------------------------%
\subsection{Função Quadrática}

É uma função polinomial de grau 2, ou seja, tem a forma
\[
  f(x)=ax^2+bx+c,
\]
com $a, b, c \in \R$ e $a \neq 0$.
O gráfico desse tipo de função é sempre uma parábola e tem concavidade para
cima se $a>0$ e concavidade para baixo se $a<0$. O gráfico de uma parábola pode
ser construído a partir de translações da parábola $ax^2$. A figura a seguir
mostra o exemplo de duas parábolas.

%------------------------------------------------------------------------------%
\begin{figure}
  \centering
  \includegraphics[scale=0.5]{p1}
  \includegraphics[scale=0.8]{p2}
  \caption{Gráfico das funções $y=x^2+x+1$ e $-2x^2+3x+1$.}
\end{figure}
%------------------------------------------------------------------------------%

%------------------------------------------------------------------------------%
\subsection{Função Raiz}

É uma função do tipo
\[
  f(x) = \sqrt[n]{x},
\]
com $n$ sendo uma constante natural maior que 1.

O domínio é $\R$ se o índice $n$ é ímpar e é $[0, \infty)$ se o
índice é par.

O conjunto imagem é $\R$ se o índice $n$ é ímpar e é
$[0, \infty)$ se o índice é par.

Há um único zero em $x=0$.

A função é crescente em todo o domínio.

Quando o índice é ímpar, $f$ é ímpar.

%------------------------------------------------------------------------------%
\begin{figure}
  \centering
  \includegraphics[scale=0.9]{raiz} \quad
  \includegraphics[scale=0.9]{raiz-cubica}
  \caption{Gráficos das funções $y=\sqrt{x}$ e $\sqrt[3]{x}$, respectivamente}
\end{figure}
%------------------------------------------------------------------------------%

%------------------------------------------------------------------------------%
\subsection{Função Recíproca}

É uma função do tipo
\[
  f(x) = \frac{1}{x^n},
\]
com $x \neq 0$.

O domínio inclui todos os valores reais, exceto o zero.

O conjunto imagem é $(-\infty,0) \cup (0, \infty)$ se o expoente é ímpar e
$(0, \infty)$ se o expoente é par.

A função não tem zeros.

Os exemplos a seguir ilustramQuando o expoente é ímpar, a função é decrescente
em todos os pontos do domínio. Já quando o expoente é par, $f$ é crescente para
$x<0$ e decrescente para $x>0$.

Quando o expoente é par, $f$ é par. Quando o expoente é ímpar, $f$ é ímpar.

%------------------------------------------------------------------------------%
\begin{figure}
  \centering
  \includegraphics[scale=0.7]{hiperbole}
  \caption{Gráfico da função $\frac{1}{x}$.}
\end{figure}
%------------------------------------------------------------------------------%

%------------------------------------------------------------------------------%
\subsection{Função Racional}

É uma função do tipo
\[
  f(x)=\frac{p(x)}{q(x)},
\]
onde $p(x),q(x)$ são polinômios e $q(x) \neq 0$.

Por exemplo, a função
\[
  f(x) = \frac{2x^4-x^2+1}{x^2-4}
\]
é uma função racional e seu
domínio é \(\left\{ x \in \R \st x \neq \pm 2 \right\}\).

%------------------------------------------------------------------------------%
\subsection{Função exponencial}

É uma função do tipo
\[
  f(x) = a^x.
\]
onde $a$ é uma constante positiva.

O domínio de $f$ é $\R$ (ou $(- \infty, \infty)$).

A imagem é $(0, \infty)$.

A função não possui zeros;

Se $a>1$, a função é crescente em todo o domínio e se $0<a<1$ a função
é decrescente em todo o domínio.

A função não é par nem ímpar.

A Figura~\ref{fig:funcao_exponencial}
mostra o gráfico das funções $y=2^x$ e $y=(0.5)^x$

\begin{figure}
  \centering
  \includegraphics[scale=0.7]{expone}
  \caption{Gráficos das funções $y=2^x$ e $y=(0.5)^x$.
    Figura retirada do livro
    Cálculo 1 do autor James Stewart.}
  \label{fig:funcao_exponencial}
\end{figure}

%------------------------------------------------------------------------------%
\subsection{Funções Trigonométricas}

As funções trigonométricas são funções angulares obtidas através do auxílio do
círculo trigonométrico. As principais funções trigonométricas são as funções
\[
  y = \sin(x), \qquad
  y = \cos(x) \quad\text{ e }\quad
  y = \tan(x).
\]
Essas funções são periódicas, ou seja, os valores da função $(f(x) = y)$ se
repetem para determinados valores da variável $x$, ou seja, para cada período
determinado pelos valores de $ x$, iremos obter valores repetidos para a
função.

O domínio de ambas funções $y=\sin(x)$ e $y=\cos(x)$ é $\R$.

A Imagem de ambas funções $y=\sin(x)$ e $y=\cos(x)$ é $[-1, 1]$, ou
seja, $-1 \leq \sin(x) \leq 1$ e $-1 \leq \cos(x) \leq 1$.

A função $y=\sin(x)$ é ímpar e a função $y= \cos(x)$ é par.

%------------------------------------------------------------------------------%
\begin{figure}
  \centering
  \includegraphics[scale=0.8]{seno}
  \includegraphics[scale=0.8]{cos}
  \caption{Gráfico da função $y=\sin(x)$ e  $y=\cos(x)$, respectivamente.}
\end{figure}
%------------------------------------------------------------------------------%

%------------------------------------------------------------------------------%
\subsection{Funções Definidas por Partes}

Uma função definida por partes é uma função que tem mais de uma lei para o seu
domínio.

A função
\[
  f(x) =
  \begin{cases}
    x^2-x,  & \text{se } x < 1    \\
    x-1,    & \text{se } x \geq 1
  \end{cases}
\]
é um exemplo de função definida por partes.

%-----------------------------------------------------------------------------%
